\chapter{System Architecture and Design}
\label{ch:system_architecture}

\section{High-Level System Architecture}
The Track platform is engineered around a modern, decoupled microservices paradigm organized as a unified monorepo. Figure \ref{fig:high_level_arch} illustrates the high-level system topology, highlighting the interaction between external clients, edge networks, compute services, and persistence layers.

\begin{figure}[htbp]
\centering
\begin{tikzpicture}[
    node distance=1.5cm and 2.0cm,
    block/.style={rectangle, draw=blue!70, fill=blue!5, thick, rounded corners, minimum width=2.8cm, minimum height=1cm, align=center, font=\small},
    service/.style={rectangle, draw=purple!70, fill=purple!5, thick, rounded corners, minimum width=2.8cm, minimum height=1cm, align=center, font=\small},
    db/.style={cylinder, draw=green!60!black, fill=green!5, thick, shape border rotate=90, minimum width=2.2cm, minimum height=1.2cm, align=center, font=\small},
    edge/.style={rectangle, draw=orange!70, fill=orange!5, thick, rounded corners, minimum width=3.2cm, minimum height=1cm, align=center, font=\small},
    line/.style={draw, -latex', thick}
]
    % Nodes
    \node [block] (client) {React 18 SPA\\(Browser Client)};
    \node [edge, right=of client] (cf) {CloudFront Edge\\(track.mourya.tech)};
    \node [block, above=1.2cm of cf] (s3) {Amazon S3\\(Static Web Bundle)};
    \node [service, right=of cf] (apigw) {HTTP API Gateway v2\\(AWS ap-south-1)};
    \node [service, right=of apigw] (lambda) {Fastify 5 Lambda\\(Node.js 20 ARM64)};
    \node [db, below=1.5cm of lambda] (neon) {Neon PostgreSQL\\(Serverless DB)};
    \node [service, below=1.5cm of apigw] (ai) {AI Microservice\\(FastAPI / PyTorch)};
    \node [block, below=1.5cm of client] (gmail) {Gmail Pub/Sub\\(OAuth2 Webhook)};
    
    % Connections
    \draw [line] (client) -- node[above, font=\footnotesize]{HTTPS} (cf);
    \draw [line] (cf) -- node[left, font=\footnotesize]{Origin} (s3);
    \draw [line] (cf) -- node[above, font=\footnotesize]{/api/*} (apigw);
    \draw [line] (apigw) -- node[above, font=\footnotesize]{Proxy} (lambda);
    \draw [line] (lambda) -- node[right, font=\footnotesize]{Drizzle ORM} (neon);
    \draw [line] (lambda) -- node[above, font=\footnotesize]{HTTP RPC} (ai);
    \draw [line] (gmail) -- node[above, font=\footnotesize]{Webhook} (lambda);
    \draw [line] (ai) |- node[below, font=\footnotesize]{Embeddings} (neon);
\end{tikzpicture}
\caption{Track High-Level System Architecture and Component Interaction}
\label{fig:high_level_arch}
\end{figure}

The architecture enforces strict separation of concerns across tiers:
\begin{enumerate}
    \item \textbf{Edge Delivery Tier}: Amazon CloudFront terminates TLS sessions using an AWS Certificate Manager (ACM) SSL certificate bound to the custom domain \texttt{track.mourya.tech}. Requests for static assets are served from Amazon S3, while dynamic API calls matching \texttt{/api/*} are forwarded directly to the API Gateway.
    \item \textbf{Application Compute Tier}: Built on Fastify 5 and bundled into a self-contained ES module executed within an AWS Lambda ARM64 runtime. Fastify handles routing, input schema validation, session cookies, and database transactions.
    \item \textbf{Persistence and Queue Tier}: Hosted on Neon Serverless PostgreSQL. In addition to relational business data, the database houses the transactional outbox queue, managing background job synchronization with row-level locks.
    \item \textbf{AI Inference Tier}: An asynchronous FastAPI service hosting PyTorch neural models (\texttt{GLiNER} and \texttt{BGE-small}) to perform embedding generation, zero-shot entity recognition, and resume optimization.
\end{enumerate}

\section{Database Schema and Entity-Relationship Model}
The persistence layer utilizes PostgreSQL 15, managed through Drizzle ORM. The relational schema is structured around normalized tables designed to maintain strict referential integrity.

\begin{figure}[htbp]
\centering
\begin{tikzpicture}[
    node distance=1.2cm and 1.8cm,
    entity/.style={rectangle, draw=black!70, fill=gray!10, thick, rounded corners, minimum width=3.2cm, minimum height=1.6cm, align=left, font=\scriptsize},
    rel/.style={draw, -latex', thick}
]
    \node [entity] (users) {
        \textbf{users} \\
        \rule{2.8cm}{0.4pt} \\
        \underline{id}: UUID (PK) \\
        email: VARCHAR \\
        password\_hash: VARCHAR \\
        google\_id: VARCHAR \\
        created\_at: TIMESTAMP
    };

    \node [entity, right=of users] (apps) {
        \textbf{applications} \\
        \rule{2.8cm}{0.4pt} \\
        \underline{id}: UUID (PK) \\
        user\_id: UUID (FK) \\
        company\_name: VARCHAR \\
        job\_title: VARCHAR \\
        current\_stage: VARCHAR \\
        is\_stage\_locked: BOOLEAN \\
        created\_at: TIMESTAMP
    };

    \node [entity, right=of apps] (events) {
        \textbf{application\_status\_events} \\
        \rule{2.8cm}{0.4pt} \\
        \underline{id}: UUID (PK) \\
        application\_id: UUID (FK) \\
        from\_stage: VARCHAR \\
        to\_stage: VARCHAR \\
        actor: VARCHAR \\
        created\_at: TIMESTAMP
    };

    \node [entity, below=of users] (pages) {
        \textbf{company\_career\_pages} \\
        \rule{2.8cm}{0.4pt} \\
        \underline{id}: UUID (PK) \\
        company\_name: VARCHAR \\
        provider: VARCHAR \\
        slug: VARCHAR \\
        last\_synced\_at: TIMESTAMP
    };

    \node [entity, below=of apps] (jobs) {
        \textbf{job\_postings} \\
        \rule{2.8cm}{0.4pt} \\
        \underline{id}: UUID (PK) \\
        career\_page\_id: UUID (FK) \\
        normalized\_company: VARCHAR \\
        job\_title: VARCHAR \\
        job\_url: VARCHAR \\
        is\_active: BOOLEAN
    };

    \node [entity, below=of events] (msgs) {
        \textbf{mailbox\_messages} \\
        \rule{2.8cm}{0.4pt} \\
        \underline{id}: UUID (PK) \\
        user\_id: UUID (FK) \\
        message\_id: VARCHAR \\
        sender: VARCHAR \\
        parsed\_intent: VARCHAR \\
        processed\_at: TIMESTAMP
    };

    % Relationships
    \draw [rel] (users) -- node[above, font=\tiny]{1:N} (apps);
    \draw [rel] (apps) -- node[above, font=\tiny]{1:N} (events);
    \draw [rel] (pages) -- node[above, font=\tiny]{1:N} (jobs);
    \draw [rel] (users) -- node[left, font=\tiny]{1:N} (msgs);
    \draw [rel] (apps) -- node[right, font=\tiny]{0:N} (msgs);
\end{tikzpicture}
\caption{Entity-Relationship Schema Diagram of Core Database Tables}
\label{fig:er_diagram}
\end{figure}

Key architectural tables include:
\begin{itemize}
    \item \texttt{applications}: Stores candidate applications with normalized company, role, stage, and the boolean \texttt{is\_stage\_locked} guard flag.
    \item \texttt{application\_status\_events}: Append-only audit table tracking every stage transition, recording the prior stage, destination stage, triggering actor (\textit{user} vs. \textit{email\_webhook}), and transition rationale.
    \item \texttt{company\_career\_pages} \& \texttt{job\_postings}: Relational catalog of monitored companies and their extracted positions with deduplication hashes.
    \item \texttt{password\_resets}: Stores hashed, single-use, 15-minute verification codes for secure password recovery.
\end{itemize}

\section{Application Lifecycle State Machine and Stage Lock}
Candidate applications traverse a formalized finite state machine (FSM). Figure \ref{fig:state_machine} diagrams the permissible transitions between stages.

\begin{figure}[htbp]
\centering
\begin{tikzpicture}[
    node distance=1.8cm,
    state/.style={circle, draw=blue!80, fill=blue!10, thick, minimum size=1.8cm, align=center, font=\footnotesize},
    final/.style={circle, draw=red!80, fill=red!10, thick, minimum size=1.8cm, align=center, font=\footnotesize},
    success/.style={circle, draw=green!70!black, fill=green!10, thick, minimum size=1.8cm, align=center, font=\footnotesize},
    trans/.style={draw, -latex', thick}
]
    \node [state] (saved) {Saved};
    \node [state, right=of saved] (applied) {Applied};
    \node [state, right=of applied] (screening) {Screening};
    \node [state, right=of screening] (interview) {Interview};
    \node [success, above right=0.8cm and 1.5cm of interview] (offer) {Offer};
    \node [final, below right=0.8cm and 1.5cm of interview] (rejected) {Rejected};

    \draw [trans] (saved) -- node[above, font=\tiny]{Submit} (applied);
    \draw [trans] (applied) -- node[above, font=\tiny]{Recruiter Reachout} (screening);
    \draw [trans] (screening) -- node[above, font=\tiny]{Technical / Onsite} (interview);
    \draw [trans] (interview) -- node[above left, font=\tiny]{Accepted} (offer);
    \draw [trans] (interview) -- node[below left, font=\tiny]{Rejection} (rejected);
    \draw [trans] (applied) to[bend right=30] node[below, font=\tiny]{Direct Rejection} (rejected);
    \draw [trans] (screening) to[bend right=20] node[below, font=\tiny]{Rejection} (rejected);
    
    % Rollback loop
    \draw [trans, dashed] (rejected) to[bend left=45] node[above, font=\tiny]{Revert Transition} (interview);
\end{tikzpicture}
\caption{Application Lifecycle State Machine Diagram with Rollback Capability}
\label{fig:state_machine}
\end{figure}

\subsection{The Stage Lock Mechanism}
A critical vulnerability in automated career trackers is the ``destructive overwrite problem'': when a user manually schedules an upcoming final-round interview (advancing status to \texttt{interviewing}), an asynchronous rejection email arriving from a previous sub-application or older requisition might trigger an automated downgrade to \texttt{rejected}.

To permanently resolve this, Track introduces the \textbf{Stage Lock} protocol:
\begin{equation}
    \Delta(\text{Stage}) = \begin{cases}
        \text{Reject Event} & \text{if } \text{Actor} = \text{AutomatedAgent} \land \text{StageLocked} = \text{True} \\
        \text{Apply Transition} & \text{otherwise}
    \end{cases}
\end{equation}
When \texttt{is\_stage\_locked} is active, automated background workers are strictly barred from mutating application stages, instead queueing a notification for user confirmation.

\section{Staged Blocking and 3-Tier Cascade Deduplication Pipeline}
Comparing all $N$ incoming job listings across all job boards pairwise exhibits quadratic time complexity $O(N^2)$. When evaluating 50,000 job postings, this requires $1.25 \times 10^9$ pairwise comparisons—rendering neural vector inference mathematically unfeasible for interactive response times.

To solve this, we engineer a **Staged Blocking + 3-Tier Cascade Deduplication Architecture** (Algorithm \ref{alg:dedup}).

\subsection{Mathematical Formulation}

\paragraph{Stage 0: Company Normalization and Blocking}
Corporate legal entities are stripped of suffixes ($\text{"Inc."}$, $\text{"LLC"}$, $\text{"Corp."}$, $\text{"Ltd."}$, $\text{"Technologies"}$) and special punctuation via regular expressions:
\begin{equation}
    \text{Norm}(C) = \text{RegexReplace}(C, \text{"\textbackslash b(inc|llc|corp|ltd|technologies)\textbackslash b"}, \text{""})
\end{equation}
Pairwise comparisons are restricted to listings sharing identical normalized company keys $\text{Norm}(C)$, reducing total pairs from $O(N^2)$ to:
\begin{equation}
    \sum_{m=1}^{M} \frac{k_m(k_m - 1)}{2} \approx O(k^2) \quad \text{where } k_m \ll N
\end{equation}

\paragraph{Tier 1: Combined Jaccard Lexical Filter}
For candidate job pair $(A, B)$ within a company block, compute the combined word token and character 3-gram Jaccard score:
\begin{equation}
    J_{\text{comb}}(A, B) = 0.6 \cdot J_{\text{token}}(A, B) + 0.4 \cdot J_{\text{char3}}(A, B)
\end{equation}
\textit{Early Exit}: If $J_{\text{comb}}(A, B) \ge 0.82$, mark as duplicate and bypass all subsequent tiers.

\paragraph{Tier 2: Token Sort Levenshtein Edit Ratio}
If $J_{\text{comb}} < 0.82$, evaluate character edit distance over alphabetically sorted tokens:
\begin{equation}
    R_{\text{sort}}(A, B) = \text{TokenSortRatio}(A, B)
\end{equation}
\textit{Early Exit}: If $R_{\text{sort}}(A, B) \ge 0.86$ or $\frac{J_{\text{comb}} + R_{\text{sort}}}{2} \ge 0.76$, mark as duplicate.

\paragraph{Tier 3: Dense Vector Semantic Cosine Evaluation}
If the lexical average falls into the ambiguous candidate zone ($0.35 \le \text{Lexical} < 0.76$), invoke \texttt{bge-small-en-v1.5} to produce 384-dimensional embeddings $\mathbf{u}, \mathbf{v}$ over $\text{Title} + \text{Snippet}$:
\begin{equation}
    S_{\text{semantic}}(A, B) = \frac{\mathbf{u} \cdot \mathbf{v}}{\|\mathbf{u}\|_2 \|\mathbf{v}\|_2}
\end{equation}
If $S_{\text{semantic}}(A, B) \ge 0.85$, confirm pair as a semantic duplicate. If $\text{Lexical} < 0.35$, immediately discard pair without neural evaluation.

\section{Security Architecture}
Security controls are enforced across two primary vectors:

\subsection{Server-Side Request Forgery (SSRF) Defense}
Because Track allows administrators to configure arbitrary career portal URLs for synchronization, malicious users could theoretically submit internal AWS loopback addresses to access the cloud metadata service (\texttt{169.254.169.254}) or internal services. 

To mitigate this, the HTTP fetch client resolves the target DNS hostname before issuing requests, validating the IP address against an unalterable blacklist blocking RFC1918 private networks, link-local subnets, and loopback ranges.

\subsection{Fernet Authenticated OAuth Token Encryption}
To store third-party Gmail OAuth2 refresh tokens securely at rest, tokens are encrypted using **Fernet** authenticated symmetric encryption. Fernet guarantees that data cannot be manipulated or read without the cryptographic key, combining:
\begin{itemize}
    \item **AES-128 in CBC Mode**: Encrypting plaintext tokens using PKCS7 padding.
    \item **HMAC-SHA256**: Computing an authentication message digest over the initialization vector and ciphertext to detect any tampering at rest.
\end{itemize}

